\BourbakiExerciseGroup

\begin{enumerate}
\def\labelenumi{\arabic{enumi})}
\item
  设 A 是一个不一定交换的环，非零且无零因子（I，第 98 页）。证明：若存在整数 \(m \geqslant 1\) 使得 \(m\mathbf{A} = 0\) ，则具有此性质的所有整数组成的集合是一个理想 \(p\mathbf{Z}\) ，其中 \(p\) 是一个素数，且 A 是特征为 \(p\) 的环。
\item
  设 \(m, n\) 为满足 \(0 \leqslant n \leqslant m\) 的整数。设 \(p\) 为一个素数，并设 \(m = \alpha_0 + \alpha_1p + \cdots + \alpha_rp^r\) 与 \(n = \beta_0 + \beta_1p + \cdots + \beta_rp^r\) 为 \(m\) 和 \(n\) 在基 \(p\) 下的展开式，因而我们有 \(0 \leqslant \alpha_i < p\) ， \(0 \leqslant \beta_i < p\) （对于 \(0 \leqslant i \leqslant r\) ）；假设 \(\alpha_r \neq 0\) 。证明 \(\binom{m}{n} \equiv 0 (\bmod p)\) ，只要 \(\beta_i > \alpha_i\) 对至少一个指标 \(i\) 成立。另一方面，如果 \(\beta_i \leqslant \alpha_i\) （对于 \(0 \leqslant i \leqslant r\) ），则
\end{enumerate}

\[
\binom {m} {n} \equiv \binom {\alpha_ {0}} {\beta_ {0}} \dots \binom {\alpha_ {r}} {\beta_ {r}} \pmod {p}  .
\]

（在 \((\mathbf{Z} / p\mathbf{Z})[\mathbf{X}]\) 中，我们有 \((1 + \mathbf{X})^m = (1 + \mathbf{X})^{\alpha_0}(1 + \mathbf{X}^p)^{\alpha_1}\dots (1 + \mathbf{X}^{p^r})^{\alpha_r})\)

\begin{enumerate}
\def\labelenumi{\arabic{enumi})}
\setcounter{enumi}{2}
\tightlist
\item
  设 \(\mathbf{A}\) 为一个交换环。对于每个多项式 \(f \in \mathbf{A}[\mathbf{X}]\) ，我们在多项式环 \(\mathbf{A}[\mathbf{X}, \mathbf{Y}]\) （\(\mathbf{A}\) 上的两个未定元的多项式环）中写下
\end{enumerate}

\[
f (\mathbf {X} + \mathbf {Y}) = \sum_ {m = 0} ^ {\infty} \Delta_ {m} f (\mathbf {X}) \mathbf {Y} ^ {m},
\]

其中 \(\Delta_{m}\) 因而是 A{[}X{]} 到自身的一个 A-线性映射。\\
a) 若 Z 是一个未定元，且 \(\tau_{Z}\) 是 A{[}X{]} 到 A{[}X, Z{]} 的一个 A-线性映射，满足 \(\tau_{Z}f = f(X + Z)\) ，证明 \(\tau_{Z} \circ \Delta_{m} = \Delta_{m} \circ \tau_{Z}\) （其中右侧的 \(\Delta_{m}\) 是上面在环 B{[}X{]} 中定义的映射，其中 B = A{[}Z{]}）。由此推出，在 A{[}X{]} 中我们有

\[
\Delta_ {m} \Delta_ {n} = \Delta_ {n} \Delta_ {m} = \binom {m + n} {m} \Delta_ {m + n}
\]

对于整数 \(m, n \geqslant 0\) （参见 IV，第 8 页）。

\begin{enumerate}
\def\labelenumi{\alph{enumi})}
\setcounter{enumi}{1}
\tightlist
\item
  设 A 是特征为 p \textgreater{} 0 的环，并且对每个整数 \(k \geqslant 0\) 令 \(D_{k} = \Delta_{p^{k}}\) 。证明：若 \(f \in A[X]\) 且 \(g \in A[X^{p^{k}}]\) ，则 \(D_{k}(fg) = g \cdot D_{k}f + f \cdot D_{k}g\) 。
\end{enumerate}

每个多项式 \(f \in \mathbf{A}[\mathbf{X}]\) 都可以唯一地写成形式 \(f(\mathbf{X}) = \sum_{m=0}^{\infty} \mathbf{X}^{mp^k} g_m(\mathbf{X})\) ，其中 \(\deg(g_m) < p^k\) ；证明

\[
\mathrm{D} _ {k} f (\mathbf {X}) = \sum_ {m = 0} ^ {\infty} m \mathbf {X} ^ {(m - 1) p ^ {k}} g _ {m} (\mathbf {X})
\]

并推导出 \(D_{k}^{p}=0\)。

\begin{enumerate}
\def\labelenumi{\alph{enumi})}
\setcounter{enumi}{2}
\tightlist
\item
  对于每个整数 m \textgreater{} 0，设 \(m = \alpha_{0} + \alpha_{1}p + \cdots + \alpha_{k}p^{k}\) 为 m 在基 p 下的展开式（集合论，III，第 177 页），因而 \(0 \leqslant \alpha_{j} < p\) （对于 \(0 \leqslant j \leqslant k\) ）；假设 \(\alpha_{k} \neq 0\) 。证明
\end{enumerate}

\[
\Delta_ {m} = \left(\frac {1}{\alpha_ {0} !} D _ {0} ^ {\alpha_ {0}}\right) \left(\frac {1}{\alpha_ {1} !} D _ {1} ^ {\alpha_ {1}}\right) \dots \left(\frac {1}{\alpha_ {k} !} D _ {k} ^ {\alpha_ {k}}\right).
\]

（如果 \(\Delta_m'\) 是该关系的右端，注意如果我们引入展开式 \(f(\mathbf{X} + \mathbf{Y}) = \sum_{q=0}^{\infty} f_q(\mathbf{Y}) \mathbf{X}^q\) ，我们有 \(\Delta_m' f(\mathbf{X} + \mathbf{Y}) = \sum_{q=0}^{\infty} f_q(\mathbf{Y}) \Delta_m'(\mathbf{X}^q)\) ，并推出它可以写成 \(\Delta_m' f(\mathbf{X} + \mathbf{Y}) = f_m(\mathbf{X}) + \mathbf{Y} g(\mathbf{X}, \mathbf{Y})\) ，其中 \(g\) 是一个多项式。）

\begin{enumerate}
\def\labelenumi{\arabic{enumi})}
\setcounter{enumi}{3}
\tightlist
\item
  设 \(\mathbf{G}\) 是一个交换群，以加法书写，并设 \(p\) 是一个素数。用 \(\mathbf{G}\left[\frac{1}{p}\right]\) 表示交换群 \(\mathbf{G} \otimes_{\mathbf{Z}} \mathbf{Z}\left[\frac{1}{p}\right]\) 。证明：若 \(\mathbf{K}\) 是一个特征为 \(p\) 的域，则群代数 \(\mathbf{K}[\mathbf{G}]\) 的完全闭包是群代数 \(\mathbf{K}^{p^{-\infty}}\left[\mathbf{G}\left[\frac{1}{p}\right]\right]\) 。
\end{enumerate}

* 5) 设 \(\mathfrak{P}\) 为素数的集合，并令 \(\mathbf{A} = \prod_{p\in \mathfrak{P}}\mathbf{F}_p\) 。对于每个滤子 \(\mathfrak{F}\) （\(\mathfrak{P}\) 上的）

（一般拓扑学，I，第 57 页）用 \(m_{\mathfrak{F}} \subset \mathbf{A}\) 表示由满足如下条件的 \(u = (u_p)_{p \in \mathfrak{P}}\) 组成的集合：\(p \in \mathfrak{P}\) 中使 \(u_p = 0\) 成立的那些 p 组成的集合是 \(\mathfrak{F}\) 的一个成员。证明 \(m_{\mathfrak{F}}\) 是一个理想，并且我们由此得到 \(\mathbf{A}\) 的异于 \(\mathbf{A}\) 的理想与 \(\mathfrak{P}\) 上的滤子之间的一个双射。设 \(\mathfrak{F}\) 是 \(\mathfrak{P}\) 上的一个在离散拓扑中不具有极限的超滤子（一般拓扑学，I，第 60 页）。证明 \(\mathbf{A}/m_{\mathfrak{F}}\) 是一个特征为 0 的域，同构于 \(\mathbf{C}\) 的一个子域。 *

